\section{Original-risk marginal acquisition, training and policy control}
\label{supp:marginal121}

This section proves the statements in Section~\ref{sec:marginal121}.
The economic assumptions are the exact primitives in
\eqref{eq:investment95}--\eqref{eq:cost95}, with \(d=2\),
\(\beta=15/16\), \(R=1/32\), and a separately fixed
\(\alpha\in\{1/2,3/4\}\).
The original strong convexity, domain invariance and constrained value
gradient are proved in Proposition~\ref{prop:investmentfoundations119};
the original \(V_1\) gradient modulus \(M_\alpha\) is proved in
Theorem~\ref{thm:modulus120}.
Neither property is an assumption about a fitted critic.
The null-space identity, finite-difference arguments, box dual bound
and constrained strong-convexity inequalities use elementary or
classical arguments; their proofs are included to specify exactly
which training and policy claims they support.

\subsection{Constrained motion of the original adverse tail}

Let
\begin{gather*}
 \boldsymbol b=(1/2,1/4)^\top,\qquad
 \boldsymbol s=(1,-1)^\top,\\
 B=2R(1-\alpha),\qquad L=2R\alpha,\qquad
 I=[-R,R-L].
\end{gather*}
For \(u=(x,v)\in\Omega\), put
\[
 m(u)=F(x,A(x)v,0),\qquad
 f_u(z)=V_1(m(u)+\boldsymbol s z),\qquad
 H_u^{\rm cen}(c)=\int_c^{c+L}f_u(z)\,dz .
\]
The letter \(H_u^{\rm cen}\) denotes the central-window integral
only; \(H_\alpha(x,v)\) continues to denote the action marginal.
The original invariant domain places every successor, and the
segments used below, in the convex state cube.
The inherited lower curvature of \(V_1\) is \(181/128\).
Consequently \(f_u\) is \(m_z=181/64\)-strongly convex in \(z\).
It is continuously differentiable, and
\[
 (H_u^{\rm cen})'(c)=f_u(c+L)-f_u(c)
\]
is strongly monotone with modulus \(m_zL>0\).
Thus \(H_u^{\rm cen}\) has a unique minimizer \(c(u)\in I\).

\begin{lemma}[An envelope and a constrained cut modulus]
\label{lem:cutmarginal121}
For every \(u\),
\begin{equation}\label{eq:tailmarginal121}
 \rho_\alpha(f_u)=\frac1B
       \left\{\int_{-R}^{c(u)}f_u(z)\,dz+
                         \int_{c(u)+L}^{R}f_u(z)\,dz\right\}.
\end{equation}
If two such functions satisfy
\(\operatorname{Lip}_z(f_u-f_{\widetilde u})\leq E\), their
minimizing cuts satisfy
\[
 |c(u)-c(\widetilde u)|\leq E/m_z.
\]
Both assertions include an endpoint minimizing cut and a tail of
zero length.
\end{lemma}

\begin{proof}
For any \(c\in I\), the indicator of the two tails divided by
\(t=1-\alpha\) is a feasible density relative to the original
uniform law.  Its objective is the total integral of \(f_u\)
minus \(H_u^{\rm cen}(c)\), divided by \(B\).
If the minimizing cut is interior, its two endpoint values agree.
Convexity then places all values inside the central interval below
those on its complement.  If the cut equals \(-R\), its right
derivative is nonnegative, so \(f_u(-R+L)\geq f_u(-R)\).
Convexity makes \(f_u\) nondecreasing to the right of \(-R+L\),
and every value in the central interval is at most
\(f_u(-R+L)\).  Thus the right tail is again a set of the largest
values.  The other endpoint is symmetric.
The tail density therefore maximizes the risk envelope.
One can verify maximality directly: if \(q_*\) has cap \(1/t\)
on these tails and zero on the center, choose a threshold \(\zeta\)
separating their values.  For any other feasible \(q\),
\[
 \mathbb E(q-q_*)f_u
   =\mathbb E(q-q_*)(f_u-\zeta)\leq0,
\]
using equal total mass and the signs on the two sets.
Strong convexity makes every level set contain at most two
points, which have zero uniform-law mass, so no atom or
fractional threshold mass changes this argument.

Write \(c=c(u)\), \(\widetilde c=c(\widetilde u)\).
Their constrained optimality inequalities are
\[
 (H_u^{\rm cen})'(c)(\widetilde c-c)\geq0,\qquad
 (H_{\widetilde u}^{\rm cen})'(\widetilde c)
                                      (c-\widetilde c)\geq0.
\]
Strong monotonicity of \((H_u^{\rm cen})'\), combined with these
two inequalities, gives
\begin{align*}
 m_zL|c-\widetilde c|^2
 &\leq
 \bigl\{(H_u^{\rm cen})'(\widetilde c)
       -(H_{\widetilde u}^{\rm cen})'(\widetilde c)\bigr\}
                                      (\widetilde c-c)\\
 &\leq LE\,|\widetilde c-c|.
\end{align*}
The last bound follows from the Lipschitz constant of the
difference across an interval of length \(L\).  Divide when the
cuts differ; the coincident-cut case already satisfies the
claim.  The variational inequalities are valid at the
endpoints, so no equation \((H_u^{\rm cen})'(c)=0\) is imposed
there.
\end{proof}

\begin{proof}[Proof of Theorem~\ref{thm:rootmarginal121}]
The original \(V_1\) is continuously differentiable with
\(M_\alpha\)-Lipschitz gradient on the cube, including its
continuous boundary extension.  Since the normalized feasible
set for \(c\) is compact and the minimizer unique, \(c(u)\) is
continuous.  This also follows quantitatively from
Lemma~\ref{lem:cutmarginal121} and the estimates below.
The objective under any fixed feasible cut is continuously
differentiable in its parameters, uniformly over this compact
cut set.  Comparing a value increment at the old and at the
new maximizing cut and using uniform first-order expansion
therefore gives the envelope derivative with the cut held
fixed:
\begin{equation}\label{eq:actualmarginal121}
 H_\alpha(u)=\frac{\beta}{B}
       \left\{\int_{-R}^{c(u)}m_u^a(z)\,dz
                +\int_{c(u)+L}^{R}m_u^a(z)\,dz\right\},
 \qquad
 m_u^a(z)=\boldsymbol b^\top
                   \nabla V_1(m(u)+\boldsymbol s z).
\end{equation}
The derivative of the normalized transition in \(v\) is
\(A(x)\boldsymbol b\); this factor cancels on dividing the
envelope derivative by \(A(x)\).
The argument does not differentiate \(c(u)\), an optimizing
last-date action, or the rounded deployed controller.
Continuity of the integrands and of the cuts proves the
continuous marginal extension at all boundaries.

First fix \(x\) and use the original action coordinate \(a\).
The original gradient bound gives
\begin{align*}
 |m_a^a(z)-m_b^a(z)|
       &\leq M_\alpha\|\boldsymbol b\|_2^2|a-b|
         =\frac{5M_\alpha}{16}|a-b|,\\
 \operatorname{Lip}_z m_a^a
       &\leq M_\alpha\|\boldsymbol b\|_2
                          \|\boldsymbol s\|_2
         =\frac{\sqrt{10}}4M_\alpha\leq C_\alpha,\\
 \operatorname{Lip}_z(f_a-f_b)
       &\leq C_\alpha|a-b|.
\end{align*}
Hence the cuts differ by at most
\(C_\alpha|a-b|/m_z\).
For a fixed integrand \(m_a^a\), changing a tail cut from
\(c\) to \(c'\) changes its normalized tail integral by at most
\[
 \frac1B\left|\int_{c'}^c
            \{m_a^a(z)-m_a^a(z+L)\}\,dz\right|
 \leq C_\alpha\frac{\alpha}{t}|c-c'|.
\]
Changing both the integrand and the cut in
\eqref{eq:actualmarginal121} therefore proves
\[
 |\partial_a R_\alpha(x,a/A(x))
       -\partial_a R_\alpha(x,b/A(x))|
 \leq\beta D_\alpha|a-b|.
\]
Here \(\partial_a\) means differentiation of the displayed
composition as a function of the original action.
Since \(h''(a)=2+48a^2\) lies between 2 and 5 on the original
action interval, the claimed Lipschitz constant of \(Q_x'\)
is \(L_\alpha=5+\beta D_\alpha\).
Convexity of \(V_1\), the affine action transition and the
monotone convex risk envelope make the continuation term
convex in \(a\).  The stage cost gives \(2\)-strong convexity.
Multiplying action increments by \(A(x)\) proves the
normalized curvature statement.

It remains to bound variation with both state and normalized
action.  The original calculation in
Theorem~\ref{thm:modulus120} gives, for fixed \(v\),
\(\|D_xm\|_2\leq3/4\).
Also \(D_vm=A(x)\boldsymbol b\), so
\(\|D_vm\|_2\leq\sqrt5/16\).
For the full three-dimensional parameter,
\[
 \|D_um\|_2\leq
       \sqrt{(3/4)^2+(\sqrt5/16)^2}
       =\frac{\sqrt{149}}{16}\leq\frac{49}{64}.
\]
The last inequality follows on squaring from
\(2384\leq2401\).  Integrating this Jacobian on a segment
in the convex parameter cube gives the same global
Lipschitz constant for \(m\).
Since \(\|\boldsymbol b\|_2=\sqrt5/4\leq9/16\),
\[
 |m_u^a(z)-m_{\widetilde u}^a(z)|
 \leq\frac{441M_\alpha}{1024}\|u-\widetilde u\|_2.
\]
The innovation derivative of \(f_u-f_{\widetilde u}\) is
bounded by
\[
 M_\alpha\|\boldsymbol s\|_2
             \|m(u)-m(\widetilde u)\|_2
 \leq\frac{147M_\alpha}{128}\|u-\widetilde u\|_2,
\]
using \(\sqrt2\leq3/2\).
Lemma~\ref{lem:cutmarginal121} bounds the change of cut by
this quantity divided by \(m_z\).  Applying the fixed-integrand
tail bound above and \(C_\alpha=4M_\alpha/5\) proves
\[
 |H_\alpha(u)-H_\alpha(\widetilde u)|
 \leq\beta\left\{\frac{441M_\alpha}{1024}
       +\frac{147M_\alpha^2\alpha}{160m_zt}\right\}
                 \|u-\widetilde u\|_2.
\]
All denominators are positive for the declared risk levels.
The ordinary-expectation case \(\alpha=0\) can omit the cut
argument, but no \(\alpha=1\) normalized-tail statement is
asserted here.
\end{proof}

\subsection{What the finite-risk oracle has to enclose}

The label theorem assumes an enclosure of the original
optimal-future objective \(Q_x(a)\).  A deployed terminal
policy value is a different oracle unless its error is
explicitly removed or enclosed.  For clarity, a sufficient
original-law construction is the following.  Partition
\([-R,R]\) into cells \(B_i=[z_i^-,z_i^+]\), with lengths
\(h_i\), midpoint \(z_i\), and probability \(p_i=h_i/(2R)\).
Suppose actual paid last-date calculations supply
\[
 l_i\leq V_1(F(x,a,z_i)),\qquad
 V_1(F(x,a,z_i^\pm))\leq u_i^\pm.
\]
Let \(Z_-\) be the discrete law taking values \(l_i\)
with masses \(p_i\), and let \(Z_+\) take values
\(u_i^-,u_i^+\), each with mass \(p_i/2\).
Then
\begin{equation}\label{eq:labelrisk121}
 S(x)+h(a)+\beta\rho_\alpha(Z_-)
 \ \leq Q_x(a)\leq\
 S(x)+h(a)+\beta\rho_\alpha(Z_+).
\end{equation}
Here the discrete laws are computational bounds for the original
continuous law, not substituted economic risk distributions.

\begin{proof}
For each real threshold \(\zeta\), the function
\((V_1(F(x,a,z))-\zeta)_+\) is convex in \(z\).
Its midpoint value minorizes its cell average and its endpoint
secant majorizes that average.  Monotonicity in the enclosed
node values consequently gives
\[
 \mathbb E(Z_--\zeta)_+
 \leq\mathbb E(V_1(F(x,a,Z))-\zeta)_+
 \leq\mathbb E(Z_+-\zeta)_+ .
\]
For completeness, for a bounded random variable \(W\),
\[
 \rho_\alpha(W)=
 \inf_{\zeta\in\mathbb R}
       \{\zeta+t^{-1}\mathbb E(W-\zeta)_+\}.
\]
To verify this identity, any feasible risk density satisfies
\(\mathbb E qW\leq\zeta+t^{-1}\mathbb E(W-\zeta)_+\).
Choose an \(\alpha\)-quantile as the threshold, put density
\(1/t\) above it and zero below it, and, if needed, assign
the remaining feasible mass on its atom.  This attains the
bound.  Boundedness ensures a quantile exists.
Apply the identity to the three hinge expectations while
retaining the same threshold in their pointwise inequalities,
then take infima.  Addition of the exact stage term proves
\eqref{eq:labelrisk121}.  No differentiability of a finite
quadrature, sort operation or last-date compiler is used.
\end{proof}

Every approximate last-date optimization, endpoint value,
finite-risk calculation and arithmetic operation must be
enclosed in \eqref{eq:labelrisk121}.  A valid analytic terminal
compiler can supply these intervals; its supporting state
slope is not thereby an exact derivative of \(V_1\).
The current finite-label construction needs only its valid
optimal-value intervals.

\subsection{Derivative probes, boundaries and their precision cost}

\begin{proof}[Proof of Proposition~\ref{prop:marginalprobe121}]
Fix \(x\), write \(q=Q_x\), \(g=q'(a)\), and
\(L_q=L_\alpha\).
For \(0\leq s\leq d\), strong convexity and derivative
Lipschitzness imply
\[
 2s\leq q'(a+s)-g\leq L_qs.
\]
Integrating on the forward interval gives
\[
 g+d\leq \frac{q(a+d)-q(a)}d
                         \leq g+L_qd/2.
\]
Thus the first interval in \eqref{eq:secantlabels121}
contains \(g\), including at the lower action boundary.
On the backward interval,
\(-L_qs\leq q'(a-s)-g\leq-2s\), which gives the second
interval, including at the upper boundary.
If both probes are feasible, then
\[
 -(L_q-2)s
 \leq q'(a+s)+q'(a-s)-2g
 \leq(L_q-2)s.
\]
Integrate from 0 to \(d\) and divide by \(2d\) to prove
the central radius \((L_q-2)d/4\).
The outward value differences contain the exact secants, so
the conclusions survive any declared valid oracle error.
Their intersection contains \(g\); subtraction of the known
stage derivative leaves the original continuation marginal.
If a floating implementation subtracts \(h'(a)\), it must
use a valid outward enclosure or its exactly represented
rational value.  The same requirement applies to dividing a
normalized derivative by the positive capacity.
\end{proof}

These inequalities quantify the precision allocation even
before an experiment.  If each value interval has width at
most \(2e\), the central derivative interval has half-width
at most
\[
 e/d+(L_\alpha-2)d/4.
\]
A smaller probe magnifies the paid value uncertainty.
With \(L_\alpha>2\), the unconstrained balancing choice for
this upper bound is \(d=2\sqrt{e/(L_\alpha-2)}\);
feasibility and a prospectively fixed probe rule still have
to be respected.  This formula does not prescribe changing
the current experiment's probe after seeing its labels.

\begin{proposition}[A limit of the specified three-probe interface]
\label{prop:probelimit121}
Fix an interior point with feasible symmetric step \(d>0\),
and a scalar class with derivative curvature between
\(\mu>0\) and \(L_q\geq\mu\).
An algorithm supplied only values at the three locations
\(-d,0,d\), relative to that point, cannot in general return
a derivative interval of width less than
\((L_q-\mu)d/6\), even with exact values.
If each of the three supplied value intervals has width
\(2e\), there are compatible functions whose derivatives
differ by \(2e/d\).
The central-interval radius
\(e/d+(L_q-\mu)d/4\) therefore has the correct order for
this particular information interface.
\end{proposition}

\begin{proof}
Put \(\kappa=(L_q+\mu)/2\) and, on \([-d,d]\), define
\[
 p(s)=\frac{s(s^2-d^2)}{d^2},\qquad
 q_\pm(s)=\frac{\kappa s^2}{2}
              \ \pm\frac{(L_q-\mu)d}{12}p(s).
\]
Their three values agree, because \(p(-d)=p(0)=p(d)=0\).
Also \(p'(0)=-1\), so their derivatives at zero differ
by \((L_q-\mu)d/6\).
Since \(p''(s)=6s/d^2\), both second derivatives lie in
\([\mu,L_q]\).  Extending \(p\) linearly outside each
endpoint with its endpoint derivative makes \(p\)
continuously differentiable; the second derivative then
vanishes outside.  The same curvature bounds hold on any
larger feasible interval.  Any valid output interval for the
common three exact observations must contain both
derivatives.

For the uncertainty assertion use
\[
 \widetilde q_\pm(s)=\frac{\kappa s^2}{2}\pm(e/d)s.
\]
At either endpoint their two values differ by \(2e\);
at zero they agree.  The common value intervals with midpoint
\(\kappa s^2/2\) and radius \(e\) contain both functions at
all three queried sites.  Their derivatives at zero differ
by \(2e/d\), and both curvatures are \(\kappa\).
The displayed upper radius follows by the secant proof with
\(\mu\) in place of 2.
\end{proof}

This is a local, three-observation oracle statement.  Other
paid values, gradients, global shape restrictions or access
to an original primitive formula can provide more information.
It is not a lower bound for all algorithms in the fixed
investment model, for the complete acquisition bundle, or
for a neural representation class.

\subsection{Finite-box training and exact optimization evidence}

Let each row index \(i\) include a site and one feasible
one-sided derivative branch.  The bounds
\(\ell_i\leq H_\alpha(\xi_i)\leq u_i\) are replicated at
that site's branches.  At an endpoint, only the inward
branch is meaningful.
The finite-dimensional feasible set is the same declared
box \(\Theta\) for the value-loss and marginal-loss fits
within a family and budget.

\begin{proof}[Proof of Proposition~\ref{prop:boxtraining121}]
For a scalar predicted derivative \(z\),
\[
 \max_{h\in[\ell_i,u_i]}|z-h|
      =\max\{z-\ell_i,u_i-z\}.
\]
Thus \(T(\theta)\) is exactly the robust maximum node error.
It is continuous as a finite maximum of affine functions.
The coefficient box is nonempty and compact, so its minimum
is attained.  In particular, \eqref{eq:robustlp121} is a
linear program in \((\theta,T)\), even when the represented
critic is not convex in its action.

For any feasible \(\theta\), each of the affine errors is
at most \(T(\theta)\).  Multiply the positive and negative
rows by \(y_i^+,y_i^-\), respectively, and add.  The
normalization of their combined mass gives
\[
 T(\theta)\geq
     -\sum_i\ell_i y_i^++\sum_i u_i y_i^-
                  +r^\top\theta.
\]
The minimum of the last linear term over the coordinate box
is the sum of the smaller endpoint products.  This proves
\(B_y\leq T_*\).
Each robust row error is nonnegative, so also \(0\leq T_*\).
Evaluating the finite maximum at the exactly checked stored
\(\widehat\theta\) gives \(T_*\leq P\) and the stated gap.
No assertion about a floating solver's dual feasibility has
entered the argument.
\end{proof}

All quantities in this certificate can be checked using
rational arithmetic for the declared features at the
dyadic acquisition sites.  A stored binary coefficient is
converted to the rational number represented by its bytes,
not to an intended decimal value.  Nonnegative approximate
dual weights may be rounded to rationals and renormalized;
the resulting weights, rather than their unrounded
counterparts, must satisfy the conditions of the proposition.
If all proposed weights are zero, another specified
nonnegative normalized choice gives a valid, potentially
weak bound.  A coefficient outside its box is not made
feasible by reporting the solver's constraint tolerance.

The certificate also enables an exact comparison of the
actual two learned readouts.  Evaluate
\(P_{\rm V}=T(\widehat\theta_{\rm V})\) and
\(P_{\rm M}=T(\widehat\theta_{\rm M})\) with the same rows
and labels.  If their checked values satisfy
\(P_{\rm M}<P_{\rm V}\), the marginal fit has strictly
reduced this particular robust node error.
Its certified distance from the best coefficient vector in
the box is at most \(P_{\rm M}-\max\{0,B_y\}\).
Neither inequality by itself orders off-design regret,
rejection probability or complete computational work.
The value-loss control is an actual bounded least-squares
fit.  The following certificate is evaluated on both stored
readouts, with the same exact design and midpoint targets.
It quantifies the value-loss optimization gap without
assuming that a floating solver attained the global minimum.

\begin{proposition}[An exact certificate for the value-loss control]
\label{prop:valueloss121}
Let the finite value design rows be \(\Phi_i\), let the
stored midpoint targets be \(y_i\), and set
\[
 F(\theta)=\sum_i(\Phi_i\theta-y_i)^2,
 \qquad F_* = \min_{\theta\in\Theta} F(\theta).
\]
For an exactly checked \(\widehat\theta\in\Theta\), define
\begin{align*}
 g&=2\sum_i\Phi_i^\top(\Phi_i\widehat\theta-y_i),\\
 L_F&=\max\left\{0,F(\widehat\theta)-g\cdot\widehat\theta
       +\sum_j\min\{\underline\theta_jg_j,
                              \overline\theta_jg_j\}\right\}.
\end{align*}
Then
\[
 L_F\leq F_*\leq F(\widehat\theta),\qquad
 0\leq F(\widehat\theta)-F_*\leq F(\widehat\theta)-L_F.
\]
These are finite-design optimization bounds.  They do not
remove the error in the numerical value labels or imply
uniform approximation or policy optimality.
\end{proposition}

\begin{proof}
Continuity of \(F\) and compactness of the nonempty box
give existence of its minimum.  Expanding each square,
for every \(\theta\in\Theta\) one obtains
\[
 F(\theta)=F(\widehat\theta)+g\cdot(\theta-\widehat\theta)
                 +\sum_i\{\Phi_i(\theta-\widehat\theta)\}^2.
\]
Drop the nonnegative last sum and minimize the linear
term separately over each box coordinate.  This gives
the expression inside the maximum in \(L_F\).
The additional lower bound zero follows from the sum
of squares.  Feasibility of \(\widehat\theta\) gives the
upper bound.  Subtraction proves the gap bound.
All identities hold for the actual rational numbers
represented by the checked design, targets and coefficients;
no differentiability of the Bellman oracle or fitted-feature
knot is involved in differentiating this coefficient loss.
\end{proof}

\subsection{Constructive ranges for all declared feature families}
\label{supp:featureranges121}

The stored level is the fitted or interpolated \(Q_x(0)\).
Accordingly the continuation anchor itself is that level
minus the known \(S(x)\); the production full action
critic adds \(h(A(x)v)\) to the stored level and to
the learned continuation difference.
This is the full-value identity used by the recursive
anchor theorem, not an assumption that the stored
anchor is exact off the acquisition grid.

For any fixed state, anchoring \(\phi(x,v)\) by
\(\phi(x,0)\) leaves its action derivative unchanged.
The denominator \(A(x)=(2+x_1+x_2)/16\) is affine and
at least \(1/8\).
On a rational box \(C\subset\Omega\), its range
\([A_-,A_+]\) is computed at state endpoints.
If \(D_v\phi\in[p_-,p_+]\), take the minimum and maximum
of the four rational products
\[
 p_-/A_-,\quad p_-/A_+,\quad p_+/A_-,\quad p_+/A_+
\]
to enclose \(D_v\phi/A(x)\).
This does not require independence of numerator and
denominator.  It is an outward rectangle bound.
Unions of boxes or polynomial pieces use the extrema of
their already valid bounds.

\paragraph{Quadratic and polynomial features.}
The anchored degree-two pool is
\(v,x_1v,x_2v,v^2\).
Its first three coefficients are signed.  The quadratic
family restricts the last coefficient to \([0,128]\);
the polynomial family permits \([-128,128]\).
The degree-three pool contains exactly
\[
 x_1^p x_2^q v^r,\qquad
 r\geq1,\quad p,q\geq0,\quad p+q+r\leq3.
\]
It has ten elements.  Each action derivative numerator is
\(r x_1^p x_2^qv^{r-1}\).
All coordinates lie in \([0,1]\), so endpoint powers and
products give exact range bounds for each nonnegative
monomial before division by \(A\).
The coefficients, not the basis monomials, can have either
sign; signed interval summation accounts for that fact.

\paragraph{ReLU and exactly matched hinges.}
In addition to \(v,x_1v,x_2v\), the pool has 36 features
\[
\begin{aligned}
 &\left(\frac{w_1(x_1-1/2)+w_2(x_2-1/2)+v-q}{4}\right)_+\\
 &\hspace{1em}-
 \left(\frac{w_1(x_1-1/2)+w_2(x_2-1/2)-q}{4}\right)_+ ,
\end{aligned}
\]
where \(w_1,w_2\in\{-1,0,1\}\) and
\(q\in\{1/8,3/8,5/8,7/8\}\).
For the first affine argument \(z\), the derivative
numerator is 0 when \(z<0\) and \(1/4\) when \(z>0\).
It has these two one-sided values at \(z=0\).
Compute the minimum and maximum of the affine \(z\) on
the box: a strictly negative box needs only 0, a strictly
positive box only \(1/4\), and a crossing or touching box
retains both whenever the corresponding side is present.
Using \([0,1/4]\) at every such boundary is conservative
and valid.  The 36 coefficients are nonnegative; the
three base coefficients remain signed.
A conventional coordinate-loop hinge with these same
arrays has identical mathematical derivatives and values.
A different arithmetic encoding can incur different
rounding or runtime, which must be recorded separately.

\paragraph{Clamped tensor cubic splines.}
The coarse univariate knot vector is clamped at 0 and 1
with no interior knots; the fine vector additionally has
one simple interior knot at \(1/2\).
There are respectively four and five cubic basis
functions per coordinate.  The tensor pool omits the
first action basis, leaving \(4\cdot4\cdot3=48\) and
\(5\cdot5\cdot4=100\) coefficients.
The omitted basis is the only one nonzero at \(v=0\);
thus these features already vanish at the anchor.
The standard degree-\(k\) basis is built by the finite
recursion
\[
 B_{i,k}(s)=
 \frac{s-t_i}{t_{i+k}-t_i}B_{i,k-1}(s)
 +\frac{t_{i+k+1}-s}{t_{i+k+1}-t_{i+1}}B_{i+1,k-1}(s),
\]
with a zero term for a zero denominator and the
half-open indicator bases at degree zero.
The continuous endpoint limits define clamped values.
On each knot interval this is a polynomial with rational
coefficients, and differentiating that polynomial gives
the required action derivative, including endpoint
limits.  Tensor products preserve this finite
piecewise-polynomial form.

\paragraph{Tensor linear hats.}
Use the declared five- or nine-point state mesh in each
state coordinate and the five-point action mesh.
The tensor hats that use the first action hat are omitted,
leaving \(5\cdot5\cdot4=100\) or
\(9\cdot9\cdot4=324\) coefficients.
Within a mesh box the state hats are nonnegative affine
functions and the action-hat derivative is its constant
one-sided slope, zero, or the opposite slope.
Their products give numerator bounds.
At an action knot both incident slopes are retained.
Crossing a state knot does not justify dropping an
incident polynomial piece.

For completeness, polynomial ranges in the spline or
tensor pieces can be constructed without a sampled
derivative test.  On a closed coordinate interval map
the variable to \([0,1]\) and write a degree-\(n\)
polynomial in its Bernstein basis
\[
 P(s)=\sum_{k=0}^n b_k{n\choose k}s^k(1-s)^{n-k}.
\]
The coefficients are obtained by finite rational basis
conversion.  The basis functions are nonnegative and
sum to one, so \(P(s)\) lies between the smallest and
largest \(b_k\).  Tensor products have the same
nonnegative partition-of-unity property and bound
multivariate polynomial pieces by their smallest and
largest tensor coefficients.  Intersecting a cell with
the finitely many knot boxes, or merely enclosing each
intersection in a knot box, gives valid ranges.
Affine hinge cuts can also be used to subdivide cells;
retaining both sides avoids that extra subdivision
cost.  Every option has an explicit finite construction,
and its actual arithmetic work belongs in acquisition
and certificate costs if executed.

These rules provide bounds for all original states,
including feature-knot states, and all feasible
one-sided action limits.  They do not assume that a
ReLU or tensor-grid derivative has a global Lipschitz
constant.

\subsection{Finite-cover and path-integral certificates}

\begin{proof}[Proof of Theorem~\ref{thm:quotientlearning121}]
Fix a nonempty cell \(C\).  For every \(u\in C\),
Theorem~\ref{thm:rootmarginal121} and the label at any
retained node \(\xi_i\) imply
\[
 \ell_i-\Lambda_\alpha r_i(C)
   \leq H_\alpha(u)\leq
       u_i+\Lambda_\alpha r_i(C).
\]
Taking the maximum of lower endpoints and the
minimum of upper endpoints yields the target range.
The independent feature range and interval linear
combination contain the fitted action marginal.
The difference of two real numbers in these two
intervals is bounded in absolute value by the two
opposite endpoint differences.  This proves the
cell error bound \(\epsilon(C)\).
Limits from every incident feature piece are included
by construction, and the target marginal is
continuous, so the same bound holds for the required
one-sided action derivatives.

For clarity, a version that links this certificate
directly to the training loss is also available.
Choose a reference derivative row \(A_i\) at a label
node, and define
\[
 \Delta_{ij}(C)\geq
   \sup_{u\in C}
      |A(x)^{-1}D_v\phi_j(u)-A_{ij}|,
 \qquad
 T_i(\theta)=\max\{A_i\theta-\ell_i,u_i-A_i\theta\}.
\]
The feature intervals just described give such a
\(\Delta_{ij}\) by two endpoint differences.
The triangle inequality proves
\begin{equation}\label{eq:nodetocell121}
 \sup_{u\in C}|\widehat H_\theta(u)-H_\alpha(u)|
 \leq T_i(\theta)+\Lambda_\alpha r_i(C)
                 +\sum_j|\theta_j|\Delta_{ij}(C).
\end{equation}
Take the smaller of all such bounds and
\eqref{eq:finitecover121} if desired.
Even an exactly minimized \(T_i\) leaves the target
coverage and feature-variation terms.
This explicitly identifies the additional information
required to move from fit nodes to a service state.

For fixed \(x\), the fitted continuation is continuous
and piecewise polynomial in \(v\); hence it is
absolutely continuous.  The true continuation is
continuously differentiable.
For all \(v\) except the finitely many feature knots,
\[
 \left|\partial_v
    \{\widehat R_\theta(x,v)-R_\alpha(x,v)\}\right|
             \leq A(x)\epsilon_x(v).
\]
The choice of cell may be made by a fixed ordering
of the finite cells, or by taking the smallest valid
bound.  It is measurable for the piecewise-algebraic
covers used here.
Integrating this inequality between \(a/A(x)\) and
\(b/A(x)\) gives
\eqref{eq:quotientcontrast121}.
The exact stage cost and the state-only anchor
cancel from the error increment.

Equation~\eqref{eq:future121} gives
\(-\delta\leq\eta_x(b)-\eta_x(a)\leq\delta\).
Adding that difference proves the implemented-policy
contrast.  For every feasible \(b\),
\[
 Q_x(a)-Q_x(b)
 \leq\widehat q_x(a)-\widehat q_x(b)+\mathcal D_x(a,b).
\]
Take the supremum in \(b\) and add
\(0\leq\eta_x(a)\leq\delta\) to obtain the original
optimum bound.  The supremum is nonnegative because
the term at \(b=a\) is zero.
Replacing it by a larger certified upper value
preserves the bound.  None of these steps uses
convexity of the fitted critic.
\end{proof}

The scalar all-action calculation in
\eqref{eq:learnedallaction121} is constructive.
For a rectangular or polyhedral finite cover and a
fixed represented state, split the feasible action
interval at the cover crossings and the finitely
many feature knots.  A piecewise constant
\(\epsilon_x\) makes \(\mathcal D_x(a,b)\) piecewise
affine in \(b\).  The fitted action objective,
including \(h\), is a polynomial of degree at most
four on each piece for every declared family.
The required supremum is attained at a piece
endpoint or at a zero of its derivative polynomial.
For rational coefficients, signs and isolating
intervals can be checked by polynomial remainder
sequences and rational interval refinement; outward
evaluation of all these intervals gives a certified
upper bound.  An alternative is interval subdivision
with valid polynomial range bounds.  Approximate
floating roots without an enclosing calculation
are not, by themselves, this certificate.

The associated action seminorm can be written as
\[
 \|e\|_{\rm act}
 =\sup_{x\in[0,1]^2}
      \operatorname*{ess\,sup}_{0<v<1}
               |A(x)^{-1}\partial_ve(x,v)|.
\]
The supremum over states is literal; replacing it
by an essential supremum over state space could
lose a particular queried state.
For the present absolutely continuous action
sections, this seminorm is zero exactly when
\(e(x,v)=b(x)\) is independent of \(v\) at each
state.  Indeed a zero derivative integrates to
a constant, and an action constant has zero
derivative.  This is the inherited action-null
quotient identity.  The new finite-cover
construction bounds an approximate, learned
member of that quotient under the original
integrated AVaR objective.

\subsection{Constrained actor error and the deployed future}

\begin{proof}[Proof of Proposition~\ref{prop:marginalpolicy121}]
For every feasible displacement \(d\), original
strong convexity gives
\[
 Q_x(a+d)\geq Q_x(a)+Q_x'(a)d+d^2.
\]
Taking the infimum over the original feasible
interval proves
\(Q_x(a)-V_0(x)\leq\mathcal K_x(a;Q_x'(a))\).
The map \(\mathcal K_x(a;g)\) is a supremum
of affine functions of \(g\), and hence convex.
Its maximum over a closed scalar interval is at
an endpoint, which proves the bound \(K_x(a)\).
Completion of the square gives the projected
minimizer in the statement.
At the lower boundary \(d\geq0\); if the full
derivative interval is nonnegative, \(gd+d^2\)
is minimized at zero.  At the upper boundary
the signs reverse.  This proves both exact
boundary claims.

Let \(a_0^*\) be the unique original optimizer.
Its constrained first-order condition gives
\(Q_x'(a_0^*)(a-a_0^*)\geq0\).
Strong convexity at this point implies
\[
 |a-a_0^*|^2\leq Q_x(a)-V_0(x)\leq K_x(a).
\]
Adding the single nonnegative allowance in
\eqref{eq:future121} proves the statement
for the implemented controller.

Here is a precise interpretation of the final
stationarity assertion.  At an interior smooth
action let \(\widehat g\) be the derivative of
the fitted action objective.  At an interior
kink it may be any supplied convex combination
of the two one-sided derivatives.
At either endpoint use the inward derivative.
Every such \(\widehat g\) differs from the
continuous true derivative by at most
\(\epsilon\), provided the required one-sided
errors have that bound.
Let
\[
 N_{[0,A(x)]}(a)
   =\{n:n(b-a)\leq0\ \hbox{for every feasible }b\}.
\]
A checked constrained stationarity residual
means that a supplied \(n\) in this cone
satisfies \(|\widehat g+n|\leq r\).
Use the interval
\([\widehat g-\epsilon,\widehat g+\epsilon]\),
or its intersection with any other valid
derivative interval.
For any \(g\) in this interval and any
feasible \(d\),
\[
 gd+d^2
  =\{g-\widehat g+\widehat g+n\}d+d^2-nd
  \geq-(\epsilon+r)|d|+d^2.
\]
The last expression is at least
\(-(\epsilon+r)^2/4\).
This proves \(K_x(a)\leq(\epsilon+r)^2/4\)
for the interval supplied by this error
bound.  It also explains why an unrelated
wider derivative interval need not attain
the same numerical bound.
\end{proof}

A global exact minimizer of the continuous
piecewise-polynomial fitted objective supplies
zero constrained residual: at an interior
minimum its left derivative is nonpositive
and its right derivative nonnegative; at an
endpoint the inward derivative has the
appropriate sign.  The same sign fact applies
to a checked local minimum.
An implemented polynomial-root enumeration
may have rounding error and should report its
actual residual or invoke the retained
original all-action certificate.
No theorem here identifies a floating root
with an exact stationary action.

For a different useful form, suppose an
implemented actor is known only to have
fitted objective suboptimality at most
\(\zeta\geq0\), and the original-action
marginal error is uniformly at most
\(\epsilon\) along all feasible actions at
this state.  Write
\(\Delta=Q_x(a)-V_0(x)\).
The path inequality and the strong-convexity
distance bound give
\[
 \Delta\leq \zeta+\epsilon|a-a_0^*|
                 \leq\zeta+\epsilon\sqrt{\Delta}.
\]
Solving this scalar inequality yields
\begin{equation}\label{eq:approxactor121}
 \Delta\leq
  \left\{\frac{\epsilon+
                     \sqrt{\epsilon^2+4\zeta}}2\right\}^{\!2}.
\end{equation}
The once-only allowance \(\delta\) gives the
implemented-policy bound.
This objective-gap form is weaker than the
checked-stationarity bound when \(\zeta=0\),
but remains valid without a fitted derivative
optimality certificate.

\subsection{Complete-value anchors in a finite recursive composition}

The current-state seminorm must not be used to discard a
state-dependent value when that value becomes a future
continuation.  This issue can be checked within the
original investment primitives, without another
experimental model.  For a fixed \(\gamma\in\mathbb R\),
put \(W_1(y)=V_1(y)+\gamma(y_1+y_2)\).
Since \(s_1+s_2=0\),
\[
 W_1(F(x,a,Z))
 =V_1(F(x,a,Z))
       +\gamma\{F_1(x,a,0)+F_2(x,a,0)\}.
\]
The second summand is constant in the innovation.
Cash invariance of the same original risk map proves
\[
 Q_x^{W_1}(a)-Q_x(a)
   =\beta\gamma\{F_1(x,a,0)+F_2(x,a,0)\}.
\]
The sum of the two action coefficients is \(3/4\).
Its derivative contribution is therefore
\(3\beta\gamma/4\).
An additive state-only level in a last-date action
critic leaves that last-date minimizer unchanged,
but exporting this level as a continuation changes
the root marginal as displayed.

\begin{theorem}[Anchor, quotient and policy errors in the same recursion]
\label{thm:recursiveanchor121}
Consider a finite horizon \(t=0,\ldots,T\) with
the same bounded primitives, original feasible
actions and original transition at every comparison.
Assume the Bellman and policy recursions are
well-defined, for example under the compact,
continuous, measurable-selector conditions of
Proposition~\ref{prop:finitefoundations119}.
For \(t<T\), let
\[
 Q_t^W(x,a)=c_t(x,a)+
       \beta_t\rho_{t,x,a}\{W(F_t(x,a,\cdot))\},
 \qquad V_t(x)=\inf_a Q_t^{V_{t+1}}(x,a),
\]
where \(0\leq\beta_t<\infty\) and each original
risk map is monotone and cash invariant.
All continuations under consideration are bounded
and belong to a class on which these recursions
are justified.

At each date choose an original feasible reference
action \(a_t^0(x)\) and write an acquired action
critic as
\(\widehat Q_t(x,a)=\widehat A_t(x)+\widehat D_t(x,a)\),
with \(\widehat D_t(x,a_t^0(x))=0\).
For its supplied continuation \(W_{t+1}\), suppose
\begin{align*}
 a_t^-(x)&\leq
   \widehat A_t(x)-Q_t^{W_{t+1}}(x,a_t^0(x))
                                          \leq a_t^+(x),\\
 b_t^-(x,a)&\leq
   \widehat D_t(x,a)
    -\{Q_t^{W_{t+1}}(x,a)-Q_t^{W_{t+1}}(x,a_t^0(x))\}
                                          \leq b_t^+(x,a).
\end{align*}
Let finite uniform bounds \(d_t^-,d_t^+\)
enclose the sums of these two errors over all
original feasible pairs.
Suppose the exported value satisfies
\[
 |W_t(x)-\inf_a\widehat Q_t(x,a)|\leq\zeta_t
\]
uniformly, where \(\zeta_t\geq0\) is checked.
Given \(E_T^-\leq W_T-V_T\leq E_T^+\),
define \(E_t^\pm\) by \eqref{eq:anchorrecursion121}.
Then \(E_t^-\leq W_t-V_t\leq E_t^+\) at every
date and state.

Let the implemented actor satisfy the local
original-action bound
\[
 Q_t^{W_{t+1}}(x,\pi_t(x))
       -\inf_aQ_t^{W_{t+1}}(x,a)\leq g_t
\]
uniformly, and let its continuation policy error
be bounded by \(0\leq J_{t+1}^\pi-V_{t+1}\leq P_{t+1}\).
Then
\begin{equation}\label{eq:recursivepolicy121}
 0\leq J_t^\pi-V_t
 \leq g_t+\beta_t(E_{t+1}^+-E_{t+1}^-)
                         +\beta_tP_{t+1}.
\end{equation}
Starting from the original terminal payoff,
\(P_T=0\), these bounds compose for the full
finite-horizon policy.
\end{theorem}

\begin{proof}
Monotonicity and cash invariance give, for every
state and feasible action,
\[
 \beta_tE_{t+1}^-
 \leq Q_t^{W_{t+1}}-Q_t^{V_{t+1}}
 \leq\beta_tE_{t+1}^+.
\]
This uses uniform scalar bounds on the entire
relevant continuation, not the removal of a
state-dependent term from an expectation.
The anchor and contrast inequalities add to
\[
 d_t^-\leq
        \widehat Q_t-Q_t^{W_{t+1}}\leq d_t^+.
\]
Combining these two inequalities, taking the
infimum over the identical original feasible
set, and including the two signs of
\(\zeta_t\) proves the value recurrence.
Backward induction from \(T\) gives its
claim at all dates.

For the actor, the upper comparison at
\(\pi_t(x)\) and the lower comparison at
the infimum give
\begin{align*}
 Q_t^{V_{t+1}}(x,\pi_t(x))
                  -\inf_aQ_t^{V_{t+1}}(x,a)
 &\leq
 Q_t^{W_{t+1}}(x,\pi_t(x))
                  -\inf_aQ_t^{W_{t+1}}(x,a)\\
 &\quad+\beta_t(E_{t+1}^+-E_{t+1}^-).
\end{align*}
The actual policy recursion at the same
action replaces \(V_{t+1}\) by
\(J_{t+1}^\pi\).  Monotonicity and cash
invariance bound that increment between
0 and \(\beta_tP_{t+1}\).
This proves the upper inequality in
\eqref{eq:recursivepolicy121}.
The lower inequality follows by induction:
the terminal policy uses the original
payoff; a feasible action cannot beat the
Bellman infimum; and the monotone
recursive risk cannot lower a
nonnegative future policy error.
The same proof applies if a nonnegative
terminal policy error has separately
been specified and bounded.
\end{proof}

The local quantities \(g_t\) can be supplied
by the path-integral or constrained
marginal certificate when its target is
the indicated \(Q_t^{W_{t+1}}\).
They cannot be transferred to a different
future continuation without the displayed
span term or a sharper justified comparison.
If \(E_{t+1}^-=E_{t+1}^+\), the
continuation error is a single global
constant, and its monetary action effect
does cancel.  A state-dependent anchor
typically has a nonzero span and does not
meet that condition.

For a noncash-invariant recursive utility,
state a different assumption explicitly.
Suppose the same original one-step
aggregation map, including its stage
arguments, is monotone and satisfies
\[
 |Q_t^W(x,a)-Q_t^{\widetilde W}(x,a)|
       \leq L_t\|W-\widetilde W\|_\infty
\]
uniformly on the relevant bounded class,
with a supplied finite \(L_t\).
Let \(d_t=\max\{|d_t^-|,|d_t^+|\}\).
Then the same addition and infimum
argument gives
\[
 \|W_t-V_t\|_\infty
      \leq d_t+\zeta_t+
                       L_t\|W_{t+1}-V_{t+1}\|_\infty.
\]
Comparing the actor and the infimum pays
two continuation errors, so
\[
 0\leq J_t^\pi-V_t
 \leq g_t+2L_t\|W_{t+1}-V_{t+1}\|_\infty
                          +L_tP_{t+1}.
\]
To verify these formulas, apply the
assumed absolute Lipschitz inequality
once at every action for the value
comparison, once at the actor and once
at the infimum for its local regret,
and once to the actual policy
continuation.  Monotonicity supplies the
nonnegative-policy-error induction.
No cash-translation identity is used.
An NDU or recursive-utility specification
with singular domains or unbounded
continuation derivative must establish
this assumption on an invariant domain
before using the result.

In the present two-date study, the
acquisition target is the original
\(V_1=T_1g\), evaluated by paid
primitive calculations.  Thus its
mathematical continuation approximation
in the root comparison is zero, while
the actual last-date policy error is
at most \(\varepsilon_1=2^{-20}\).
Equation~\eqref{eq:recursivepolicy121}
reduces to the root local action gap
plus \(\beta\varepsilon_1\).
Numerical uncertainty in evaluating
\(T_1g\) is already included in the
finite label and policy intervals;
it is not permission to treat that
operator evaluation as a free oracle.

\subsection{A total original-policy service and its uniform range}
\label{supp:totalservice121}

The following argument verifies the constants of the retained
exact finite-cover fallback \texttt{rational96.py}.
It does not require a favorable first proposal or convergence
of a floating polynomial-root solver.
The declared service retains a recoverable numerical failure
in the proposal or certification stages and then invokes this
fallback at the same state.  Source-identity, domain,
corrupt-model and infrastructure failures remain failed
executions; they cannot be removed from or replaced within
the evaluation sample.

\begin{proposition}[Finite all-state qualification and primitive range]
\label{prop:totalservice121}
At every represented original state \(x\in[0,1]^2\)
and either declared risk level, the exact fallback with
1025 equally spaced feasible actions, 256 innovation
cells and exact last-date value gap at most \(2^{-30}\)
returns a feasible action whose original complete-policy
regret is at most \(\tau=1/256\).
Its policy-cost interval has width at most \(1/8192\).
Every exact last-date bisection needs at most 31
iterations, within its stated cap of 80.
These are finite arithmetic guarantees; all their
actual cost is charged when entered.

The original primitive ranges satisfy
\begin{equation}\label{eq:primitiverange121}
 0\leq S\leq49/32,\qquad
 0\leq g\leq37/16,\qquad
 0\leq h\leq5/64,
\end{equation}
and hence
\[
 0\leq V_1\leq967/256,\qquad
 0\leq V_0\leq K:=21097/4096.
\]
Any two total qualified services on this same
initial-state domain have an original-policy
cost contrast in \([-\tau,\tau]\).
\end{proposition}

\begin{proof}
For \(d=2\), the target part of \(S\) is at most
\(2(5/8)^2=25/32\), its cyclic term at most
\(1/4\), and its squared hinge at most \(1/2\).
The extra target term in \(g\) is at most
\(25/32\).  The action interval lies in
\([0,1/4]\), so \(h(a)\leq1/16+1/64=5/64\).
These prove \eqref{eq:primitiverange121}.
Monotonicity and constant preservation of the
original AVaR give
\[
 V_1\leq49/32+5/64+(15/16)(37/16)=967/256,
\]
and a second application gives
\[
 V_0\leq49/32+5/64+(15/16)(967/256)
                                      =21097/4096.
\]
Nonnegativity follows from the same property
and the nonnegative primitives.

We next verify the finite terminal solve.
For a fixed terminal-decision state \(y\),
the original one-date objective has derivative
bounded in absolute value by
\[
 |Q_1'(y,a)|\leq3/4+\beta(5)(3/4)
                        =273/64<7,
\]
because each coordinate of \(\nabla g\) has
absolute value at most 5 and
\(b_1+b_2=3/4\).
The derivative is Lipschitz with constant less
than 10.  To see the latter directly, put
\(m=F(y,a,0)\), \(k=5(m_1-m_2)\).
The exact original risk representation is
\[
 Q_1(y,a)=S(y)+h(a)+
             \beta\{g(m)+\mathcal R_\alpha(k)\}.
\]
Here \(\nabla^2g\preceq13I\) almost everywhere,
\(m_a=\boldsymbol b\),
\(k_a=5/4\), and
\(\mathcal R_\alpha'\) is
\(\alpha/(10t)\)-Lipschitz.
Consequently a derivative Lipschitz constant is
\[
 5+\beta\{65/16+5\alpha/(32t)\}<10
\]
for both declared risks.  The first derivatives
match at the squared-hinge and tail-clipping
junctions, so this almost-everywhere bound
integrates to the global one.

The exact bisection starts with \([0,A(y)]\)
and evaluates its midpoint \(a_k\) after
\(k\) halvings.  The constrained optimum
remains in the bracket: a negative derivative
discards the left half and a nonnegative
derivative the right half.
Its bracket length is \(A(y)2^{-k}\).
At each step the strong-convexity lower
support has gap
\[
 \Gamma_k=-\min_{d\in[-a_k,A(y)-a_k]}
                    \{Q_1'(y,a_k)d+d^2\}.
\]
If the optimum is interior, its derivative
is zero, so
\[
 |Q_1'(y,a_k)|\leq
           10A(y)2^{-k-1},\qquad
 \Gamma_k\leq\frac{25}{64}4^{-k}.
\]
This is below \(2^{-30}\) already at \(k=15\).
If the optimum is the lower boundary,
strong convexity implies \(Q_1'(y,a_k)\geq2a_k\);
the support minimizer is then the lower
boundary and
\(\Gamma_k\leq7a_k\leq7\,2^{-k-3}\).
The upper-boundary argument is symmetric.
This is below \(2^{-30}\) at \(k=30\).
Thus at most 31 midpoint evaluations suffice,
including endpoints with zero limiting
derivative.  The exact clipped-quadratic risk,
its derivative and every arithmetic
denominator in this calculation are finite
for the declared rational inputs.

For the root, the inherited state derivative
bound \(|\partial_jV_1|\leq15/2\) gives
\[
 |Q_x'(a)|
 \leq3/4+\beta(15/2)(3/4)
                         =771/128<7.
\]
Let \(m=1024\) and use actions
\(a_j=A(x)j/m\), \(j=0,\ldots,m\).
Every feasible action is within
\(A(x)/(2m)\leq1/(8m)\) of this grid.
If \([L_j,U_j]\) is the complete-policy
interval at \(a_j\), its lower endpoint
also minorizes \(Q_x(a_j)\).
The retained convex-order width bound,
with \(K_z=15\), \(R=1/32\), \(t\geq1/4\),
terminal point error \(e=2^{-30}\), and
terminal-policy allowance \(\varepsilon_1=2^{-20}\),
gives
\begin{align}
 U_j-L_j&\leq W,\nonumber\\
 W&=\beta\left\{\frac{15}{8\cdot256^2}
                      +2e+\varepsilon_1\right\}
   =\frac{238095}{2^{33}}<2^{-15}.
 \label{eq:fallbackwidth121}
\end{align}
The convex-order inequality acts on the
original \(V_1\).  Adding
\(\beta\varepsilon_1\) to the upper endpoint
controls the actual future policy without
assuming that policy's value is convex.

The action-cover lower optimum
\[
 B_0=\min_j L_j-\frac7{8m}
\]
is valid.  Indeed each arbitrary feasible
action is within the stated distance of a
grid point, and its \(Q_x\) value is at
least that grid value minus \(7/(8m)\).
Choose \(j_*\) with the smallest \(U_j\),
with the prescribed deterministic tie rule.
Since
\(\min_jU_j-\min_jL_j\leq\max_j(U_j-L_j)\),
\[
 U_{j_*}-B_0\leq W+\frac7{8m}.
\]
The exact service rounds \(a_{j_*}\)
down to 40 binary places.  The action
remains feasible and its displacement is
less than \(2^{-40}\).
The original \(Q_x\) value changes by
less than \(7\,2^{-40}\).
Apply this perturbation to the lower
and upper endpoints while retaining
the single future allowance already
in \(U_{j_*}\).
This step uses Lipschitzness of \(Q_x\),
not of the possibly discontinuous
deployed policy value.
The exact returned width and gap are
therefore at most
\[
 W+14\,2^{-40},\qquad
 W+\frac7{8192}+7\,2^{-40},
\]
respectively.

The final floating conversion in the
retained service rounds each rational
endpoint outward by an adjacent binary64
number.  Every such endpoint has
absolute value below 8, by the
primitive range just proved and the
displayed allowances.
A nearest binary64 conversion followed
by one outward adjacent value has
absolute error less than \(2^{-48}\)
on this range.
The two endpoint conversions and the
outward final subtraction enlarge a
gap or width by less than \(2^{-46}\).
Direct rational inequalities now give
\begin{align*}
 W+14\,2^{-40}+2^{-46}&<2^{-14}<1/8192,\\
 W+7/8192+7\,2^{-40}+2^{-46}&<2^{-10}<1/256.
\end{align*}
Thus the strict margins survive the
declared conversion.
All 1025 actions use finitely many
primitive evaluations: 513 terminal
contexts per action, each with the
finite bisection bound above.
The earlier numerical stages have
finite caps and can terminate in this
fallback, so they cannot remove a
valid state from the mathematical
service's domain.

Finally each qualified service has
\(0\leq J_x-V_0(x)\leq\tau\).
Subtracting two such inequalities
with the same original optimum
gives \([-\tau,\tau]\) for their
contrast at every state.
\end{proof}

